\documentclass[12pt,letterpaper]{article}
\usepackage[utf8]{inputenc}
\usepackage[spanish]{babel}
\usepackage{geometry}
\usepackage{array}
\usepackage{tabularx}
\usepackage{graphicx}
\usepackage{enumitem}

\geometry{margin=2.5cm}

\begin{document}

\begin{center}
\textbf{Universidad Nacional Abierta y a Distancia}\\
\textbf{Vicerrectoría Académica y de Investigación}\\
\textbf{Curso: Cálculo Diferencial}\\
\textbf{Código: 100410}\\
\vspace{0.5cm}
\Large\textbf{Guía de aprendizaje -- Tarea 3 Límites y Continuidad}
\end{center}

\section{Datos de la/el Tarea}

\subsection*{Tabla 1. Tabla de descripción}

\begin{tabularx}{\textwidth}{|>{\bfseries}p{5cm}|X|}
\hline
\multicolumn{1}{|c|}{\textbf{Aspecto}} & \multicolumn{1}{c|}{\textbf{Descripción}} \\
\hline
1. Tipo de actividad & Independiente \\
\hline
2. Momento de la evaluación & Intermedio \\
\hline
3. Unidad gestora & Escuela de Ciencias Básicas Tecnología e Ingeniería ECBTI \\
\hline
4. Puntaje de la/el Tarea & 125 puntos \\
\hline
5. La actividad inicia el: & martes, 30 de septiembre de 2025 \\
\hline
6. La actividad finaliza el: & lunes, 27 de octubre de 2025 \\
\hline
7. Horas de trabajo independiente del estudiante & 24 horas \\
\hline
\end{tabularx}

\section{Descripción detallada de la actividad de aprendizaje}

Con el desarrollo de esta actividad se espera que se alcance el siguiente resultado de aprendizaje:

\textbf{Identificar los conceptos de límite y continuidad, empleando sus propiedades y métodos con precisión, en la resolución de problemas aplicados y contextualizados.}

\subsection*{La actividad consiste en:}

Revisar los contenidos del Resultado de Aprendizaje 2 en el entorno de aprendizaje y consultar el Micrositio y los Applets de GeoGebra en el Anexo 2. Seleccionar y desarrollar cinco ejercicios del Anexo 2, usando el editor de ecuaciones en Word, y publicar avances en el foro semanalmente. Finalmente, sustentar el ejercicio 5 a través de un video, subirlo a una plataforma pública, y entregar el documento con los ejercicios resueltos y el enlace al video.

\subsection*{Para el desarrollo de esta actividad se requieren los siguientes materiales y recursos:}

\begin{itemize}
\item Lectura de la guía de aprendizaje de la Tarea 3.
\item Anexo 2 -- Ejercicios Tarea 3.
\item Applets de Geogebra para la Tarea 3.
\item Micrositio del curso y Canal de YouTube.
\item Contenidos y referentes bibliográficos -- Tarea 3.
\item Visualización de los enlaces de grabación de las web conferencias y cipas dinamizados por los docentes.
\end{itemize}

\subsection*{Para el desarrollo de esta actividad debe seguir los siguientes pasos:}

Para desarrollar los ejercicios de esta actividad, es fundamental acceder a la pestaña 4 -- Límites y Continuidad, en el Micrositio de la red de curso y a nuestro canal de YouTube, donde encontrará toda la información necesaria para cada tema, junto con algunos videos orientadores. Asimismo, para fortalecer su proceso de aprendizaje, la red de curso pone a su disposición unos Applets de GeoGebra que le permitirán consolidar conceptos de manera práctica y sencilla. Recuerde revisar los enlaces del Micrositio y de los Applets de GeoGebra en el Anexo 2 - Ejercicios Tarea 3.

\subsubsection*{Paso 1:}

Revisar en el entorno de aprendizaje y evaluación las referencias bibliográficas necesarias para el Resultado de Aprendizaje 2 y participar activamente en el foro de discusión específico para la Tarea 3. Leer y analizar de manera individual los temas del entorno de aprendizaje, relacionados con el Resultado de Aprendizaje 2, como límites y sus propiedades, límites laterales, límites indeterminados, límites al infinito, límites trigonométricos, continuidad y su uso en Geogebra.

\subsubsection*{Paso 2:}

Revisar la guía de aprendizaje y la rúbrica de evaluación para entender la dinámica de la actividad. Leer el Anexo 2 -- Ejercicios Tarea 3 y elegir los ejercicios a desarrollar y anunciando la selección en el foro.

\subsubsection*{Paso 3:}

Seleccionar una serie de ejercicios A, B, C, D, o, E, y desarrollar ese mismo literal en los 5 grupos de ejercicios. La elección de ejercicios no se puede repetir, pues cada estudiante elige los que estén disponibles para desarrollar y que no hayan sido seleccionados por otro compañero en el foro. El estudiante deberá indicar en el foro el ejercicio a desarrollar, ejemplo: \textit{Andrés Pérez -- Desarrollaré el grupo de ejercicios de la letra ``A''}

\subsubsection*{Paso 4:}

Desarrollar los ejercicios seleccionados en documento Word o pdf, utilizando el editor de ecuaciones y verificar los resultados en GeoGebra. Compartir los avances semanalmente en el foro, detallando el procedimiento. Revisar regularmente el foro para recibir retroalimentación del tutor, aclarar dudas y aplicar correcciones sugeridas. Cada aporte debe ser original del estudiante.

\subsubsection*{Paso 5:}

Compilar los ejercicios en un documento Word o pdf que debe contar con los siguientes elementos:

\begin{itemize}
\item Portada.
\item Introducción.
\item Desarrollo de los cinco ejercicios seleccionados.
\item Enlace del video del ejercicio 5 (Aplicaciones) sustentado.
\item Referencias Bibliográficas en normas APA -- versión 7
\end{itemize}

\subsubsection*{Paso 6:}

\textbf{La sustentación del video debe cumplir los siguientes parámetros:}

Grabar un video explicando paso a paso la solución del ejercicio. Puede usar la cámara de un dispositivo o compartir pantalla. La sustentación debe incluir el enunciado del ejercicio, el método de resolución, cada paso explicado y el resultado final, mostrando claridad en el proceso.

\section{Indicaciones para el desarrollo y entrega de las evidencias de aprendizaje.}

Las evidencias de aprendizaje son las acciones, productos o procesos observables que se realizan y/o entregan para manifestar las capacidades, habilidades, aptitudes y actitudes adquiridas, y que, a su vez, servirán al docente para verificar y evaluar su desempeño.

\subsection*{Las evidencias a desarrollar independientemente son:}

\begin{itemize}
\item Participar en el foro de la actividad de la Tarea 3 con los avances del desarrollo de sus ejercicios en documento Word o pdf.
\item Realizar el video de la sustentación del ejercicio relacionado con las aplicaciones.
\item Entregar del trabajo compilado en el entorno de aprendizaje y evaluación, en el espacio Tarea 3 -- límites y Continuidad - Rúbrica de evaluación y entrega de la actividad.
\end{itemize}

\subsection*{Las evidencias a desarrollar colaborativamente son:}

\begin{itemize}
\item En esta actividad no se requieren evidencias de trabajo grupal.
\end{itemize}

\subsection*{Para su desarrollo y entrega tenga en cuenta las siguientes orientaciones:}

\begin{itemize}
\item Visualizar la agenda del curso para verificar los tiempos del desarrollo de la actividad, así como también, de realimentación y calificación de la misma.
\item Escoger una asignación de ejercicios del ``Anexo 2 -- Ejercicios Tarea 3'' y mencionarlo en el foro de la actividad de la Tarea 3.
\item Desarrollar individualmente los ejercicios de acuerdo con lo que se solicita en cada enunciado, aplicando correctamente el lenguaje matemático a través del editor de ecuaciones de Word y realizando por lo menos uno o dos aportes significativos por semana en el foro de la actividad.
\item Realizar la sustentación del ejercicio de aplicación a través de un video y lo deberá alojar en una plataforma de videos por demanda (Vimeo, YouTube, Loom, Teams etc.), entregando el enlace de acceso público en el documento final.
\end{itemize}

Tenga en cuenta que todos los productos escritos independientes o grupales deben cumplir con las normas de ortografía y con las condiciones de presentación que se hayan definido.

En cuanto al uso de referencias considere que el producto de esta actividad debe cumplir con las normas Elija un elemento.

En cualquier caso, cumpla con las normas de referenciación y evite el plagio académico, para ello puede apoyarse revisando sus productos escritos mediante la herramienta Turnitin que encuentra en el campus virtual.

\section{Situaciones de orden académico}

Considere que en el acuerdo 029 del 13 de diciembre de 2013, artículo 99, se considera como faltas que atentan contra el orden académico, entre otras, las siguientes: literal e) ``El plagiar, es decir, presentar como de su propia autoría la totalidad o parte de una obra, trabajo, documento o invención realizado por otra persona. Implica también el uso de citas o referencias faltas, o proponer citad donde no haya coincidencia entre ella y la referencia'' y liberal f) ``El reproducir, o copiar con fines de lucro, materiales educativos o resultados de productos de investigación, que cuentan con derechos intelectuales reservados para la Universidad''

Las sanciones académicas a las que se enfrentará el estudiante son las siguientes:

\textbf{a)} En los casos de fraude académico demostrado en el trabajo académico o evaluación respectiva, la calificación que se impondrá será de cero puntos sin perjuicio de la sanción disciplinaria correspondiente.

\textbf{b)} En los casos relacionados con plagio demostrado en el trabajo académico cualquiera sea su naturaleza, la calificación que se impondrá será de cero puntos, sin perjuicio de la sanción disciplinaria correspondiente.

\end{document}
